\documentclass[10pt,a4paper]{amsart}
\usepackage[margin=19mm]{geometry}
\usepackage{amsmath,amssymb,mathtools}
\usepackage[scr=rsfs,cal=euler]{mathalfa}
\usepackage{xfrac}
\usepackage[unicode,hidelinks,bookmarks=false]{hyperref}
\newcommand{\ie}{i.e.\ }
\newcommand{\cf}{cf.\ }
\newcommand{\resp}{respectively\ }
\newcommand{\Subsection}{\subsection}
\providecommand{\nobreakdash}{\nobreak\hbox{-}}
\setlength{\emergencystretch}{2em}
\multlinegap0pt
\usepackage{bm}
\usepackage{bbm}
\usepackage{dsfont}
\usepackage{xspace}
\usepackage{cancel}
\usepackage{mathtools}
\usepackage{amsthm}
\makeatletter
\@ifundefined{theorem}{\newtheorem{theorem}{Theorem}}{}
\@ifundefined{lemma}{\newtheorem{lemma}[theorem]{Lemma}}{}
\makeatother
\newcommand{\FRT}{\mathrm{FRT}}
\newcommand{\WPHP}{\mathrm{WPHP}}
\title[Finite Combinatorics and Fragments of Arithmetic]{Finite Combinatorics and Fragments of Arithmetic}
\author{Wei Wang}
\date{Source: arXiv:2506.17943v1 (CC BY 4.0)}
\begin{document}
\maketitle

\noindent\emph{Source and licence.} Finite Combinatorics and Fragments of Arithmetic. Version arXiv:2506.17943v1 (CC BY 4.0), distributed under the Creative Commons Attribution 4.0 International licence, as recorded in the arXiv licence metadata for this exact version and shown on the abstract page https://arxiv.org/abs/2506.17943v1. Full rights notice: licenses/arxiv-2506.17943-CC-BY-4.0.txt.\par\smallskip

\noindent\textit{Editorial scope.} Counted unit: the Theorem whose source label is \texttt{thm:FRT-WPHP} in the article (source0.tex lines 599--606, source label \texttt{thm:FRT-WPHP}), delivered together with the article's own complete original proof of it (source0.tex lines 700--769, one \texttt{proof} environment of 70 lines). No mathematics is written, repaired or re-derived: statement and proof are the article's own text line for line. The proof is detached from the statement in the source: the article states the result at lines 599--606 and prints its proof at lines 700--769, and the proof environment's own header names the statement, so the association is the article's own. Required context retained uncounted: lem:FRT-WPHP (lines 620--636). Every definition, display and label of those lines is retained unchanged. Omitted from this package and not redistributed: 1--598, 607--619, 637--699, 770--936. Nothing inside a retained excerpt is omitted or altered: every retained body is delivered exactly as the article's lines are written, so no omission receipt of any kind accompanies this package. Citations: the retained excerpts carry no citation command at all, so no bibliography entry of the article is transcribed and no reference is redistributed. Reference adaptation: none was needed and none was made; every reference inside the delivered unit resolves inside it, so no unresolved reference or citation is produced. Printed slips kept literal: no printed word, symbol, hypothesis or number of the counted statement or of its proof was altered. Layout and class: the article's own class is \texttt{amsart} with packages including amssymb, amsmath; the wrapper is the lane's pinned \texttt{amsart} editorial wrapper at 10pt on A4, which restates the theorem-like environments the retained text needs and adds the article's own macro definitions that the retained text needs (\texttt{\textbackslash FRT}, \texttt{\textbackslash WPHP}), with extra wrapper packages (bm, bbm, dsfont, xspace, cancel, mathtools). The delivered numbering is the unit's own. Assets: no retained span contains a figure environment, a graphics-inclusion command, a PDF-page inclusion, a table or a TikZ picture; no image file is copied into this package, the package declares no asset dependency and no image file is redistributed with it. Licence: version arXiv:2506.17943v1 of this article is distributed under the Creative Commons Attribution 4.0 International licence, as recorded in the arXiv licence metadata for this exact version and shown on the abstract page https://arxiv.org/abs/2506.17943v1. A full rights notice is delivered at licenses/arxiv-2506.17943-CC-BY-4.0.txt. Characters: every retained excerpt is pure ASCII, so the delivered engine, XeLaTeX with its default Latin Modern text font, reports no missing character.
\begin{lemma}\label{lem:FRT-WPHP}
Suppose that 
\begin{itemize}
 \item[(a)] $M$ is a countable model of $I\Sigma_n$ with $n > 0$,
 \item[(b)] $a,b,c,e,x$ are elements of $M$, $b > \mathbb{N}$, $k = \max\{a,c\}$,
 \item[(c)] $\varphi$ is a $\Sigma_n^M$-formula defining $C: [y]^e \times M \to c$ 
 with $y = r^{(b)}(x,e,k)$, and
 \item[(d)] $\bar{C}(\vec{d}) = \lim_s C(\vec{d},s)$ exists for all $\vec{d} \in [y]^e$.
\end{itemize}
Then either one of the followings holds,
\begin{enumerate}
 \item[(1)] There exist $s \in M$, $H \in [y]^x \cap M$ and $i < c$, s.t.
 $\bar{C}(\vec{d}) = C(\vec{d},t) = i$ for all $\vec{d} \in [H]^e$ and $t > s$.
 \item[(2)] There exists $N$, s.t. $M \preceq_{cf,\Sigma_{n+1}} N$, 
 $[0,a]^N = [0,a]^M$, $\lim_s C^N(\cdot,s)$ is undefined on some $\vec{d} \in [y]^e \cap N$ where $C^N$ is the map defined by $\varphi$ in $N$.
\end{enumerate}
\end{lemma}

\begin{theorem}\label{thm:FRT-WPHP}
For every countable $M \models I\Sigma_n$ with $n > 0$ and any $a \in M$,
there exists $N$ s.t. $M \preceq_{cf,\Sigma_{n+1}} N$, $[0,a]^M = [0,a]^N$ and
$N \models I\Sigma_n + \FRT(\Sigma_{n+1})$.

Thus,
$I\Sigma_n + \FRT(\Sigma_{n+1}) \not\vdash \WPHP(\Sigma_{n+1})$.
\end{theorem}

\begin{proof}[Proof of Theorem \ref{thm:FRT-WPHP}]
Fix a countable $M \models I\Sigma_n$ with $n > 0$ and $a \in M$.
Also fix some $b \in M - \mathbb{N}$.
By Lemma \ref{lem:FRT-WPHP}, we can build a sequence $(M_i: i \in \mathbb{N})$ s.t.
\begin{itemize}
 \item $M_0 = M$;
 \item $M_i \preceq_{cf,\Sigma_{n+1}} M_{i+1}$;
 \item $[0,a]^M = [0,a]^{M_i}$;
 \item If $x$ and $c$ are elements of $M_i$, 
 $y = r^{(b)}(x,e,k)$ with $k = \max\{a,c\}$ and
 $\varphi$ is a $\Sigma_n^{M_i}$-formula defining a function $C: [y]^e \times M \to c$ 
 s.t. $\bar{C}(\vec{d}) = \lim_s C(\vec{d},s)$ exists for all $\vec{d} \in [y]^e$,
 then one of the followings holds
 \begin{itemize}
  \item[(i)] there exist an $M_i$-finite $H$ and $s \in M_i$ s.t. $|H| \geq x$,
  \[
    M_i \models \forall t > s,\vec{d} \in [H]^e (C(\vec{d},t) = C(\vec{d},s)) \wedge H \text{ is homogeneous for } \bar{C},
  \]
  \item[(ii)] there exists $j > i$ s.t.
  \[
    M_j \models \exists \vec{d} \in [y]^e \forall s \exists t (t > s \wedge C^{M_j}(\vec{d},t) \neq C^{M_j}(\vec{d},s)),
  \]
  where $C^{M_j}$ is the function defined by $\varphi$ in $M_j$.
 \end{itemize}
\end{itemize}
Let $N = \bigcup_{i \in \mathbb{N}} M_i$.
Then $N \models I\Sigma_n$, $N$ is a $\Sigma_{n+1}$-elementary cofinal extension of $M$ and $[0,a]^N = [0,a]^M$.

For $x,e,c \in N$, let $y = r^{(b)}(x,e,k)$ where $k = \max\{a,c\}$.
Suppose that $\bar{C}: [y]^e \to c$ is $\Sigma_{n+1}^N$.
By Limit Lemma, there exists $C: [y]^e \times N \to c$
 s.t. $C$ is defined by a $\Sigma_n^N$-formula $\varphi$ and $\bar{C}(\vec{d}) = \lim_s C(\vec{d},s)$ for all $\vec{d} \in [y]^e$.
Let $i$ be s.t. $M_i$ contains $x,e,c$ and all the parameters of $\varphi$.
Then $M_i$ also contains $k$ and $y$.
As $M_i \preceq_{\Sigma_{n+1}} N$, in $M_i$, $\varphi$ also defines a function $C^{M_i}: [y]^e \times M_i \to c$,
and $\bar{C}^{M_i}(\vec{d}) = \lim_s C^{M_i}(\vec{d},s)$ exists for all $\vec{d} \in [y]^e \cap M_i$.
There are two cases.

\emph{Case 1}, (i) holds for $M_i$ and $C^{M_i}$.
Then there exists an $M_i$-finite $H$ and $s \in M_i$, s.t. $|H| \geq x$ and
\[
  M_i \models \forall t > s,\vec{d} \in [H]^e (C^{M_i}(\vec{d},t) = C^{M_i}(\vec{d},s)) \wedge H \text{ is homogeneous for } \bar{C}^{M_i}.
\]
As $M_i \preceq_{\Sigma_{n+1}} N$,
\[
  N \models \forall t > s,\vec{d} \in [H]^e \big( C(\vec{d},t) = C(\vec{d},s)) \big).
\]
It follows that $H$ is $\bar{C}$-homogeneous in $N$.

\emph{Case 2}, (ii) holds for $M_i$ and $C^{M_i}$.
Then there exist $j > i$ and $\vec{d} \in [y]^e \cap M_j$ s.t.
\[
  M_j \models \forall s \exists t( C^{M_j}(\vec{d},t) \neq C^{M_j}(\vec{d},s) ),
\]
where $C^{M_j}$ is the function defined by $\varphi$ in $M_j$.
By $M_j \preceq_{\Sigma_{n+1}} N$,
\[
  N \models \forall s \exists t( C(\vec{d},t) \neq C(\vec{d},s) ),
\]
contradicting the assumption on $C$.
This shows that $N \models \FRT(\Sigma_{n+1})$.

Finally, if in the model $M$ above there exists a $\Sigma_{n+1}^M$-injection $F: 2a \to a$,
then $[0,2a]^N = [0,2a]^M$ as $[0,a]^N = [0,a]^M$,
and the $\Sigma_{n+1}^M$-formula defining $F$ in $M$ also defines the same injection in $N$ by $M \preceq_{\Sigma_{n+1}} N$.
Hence
\[
  N \models I\Sigma_n + \FRT(\Sigma_{n+1}) + \neg \WPHP(\Sigma_{n+1}).
\]
\end{proof}

\end{document}
